\documentclass[8pt,a4paper,twocolumn]{article}
\usepackage[utf8]{inputenc}
\usepackage[spanish]{babel}
\usepackage[margin=0.6cm]{geometry}
\usepackage{listings}
\usepackage{xcolor}
\usepackage{enumitem}

% Configuración de colores para código
\definecolor{codegreen}{rgb}{0,0.6,0}
\definecolor{codegray}{rgb}{0.5,0.5,0.5}
\definecolor{codepurple}{rgb}{0.58,0,0.82}
\definecolor{backcolour}{rgb}{0.95,0.95,0.92}

\lstdefinestyle{sqlstyle}{
    backgroundcolor=\color{backcolour},
    commentstyle=\color{codegreen},
    keywordstyle=\color{blue},
    numberstyle=\tiny\color{codegray},
    stringstyle=\color{codepurple},
    basicstyle=\ttfamily\tiny,
    breakatwhitespace=false,
    breaklines=true,
    captionpos=b,
    keepspaces=true,
    showspaces=false,
    showstringspaces=false,
    showtabs=false,
    tabsize=2
}

\lstset{
    basicstyle=\ttfamily\tiny,
    keywordstyle=\color{blue}\bfseries,
    commentstyle=\color{green},
    stringstyle=\color{red},
    showstringspaces=false,
    breaklines=true,
    frame=single,
    numbers=left,
    numberstyle=\tiny\color{gray},
    inputencoding=utf8,
    literate={á}{{\'a}}1 {é}{{\'e}}1 {í}{{\'\i}}1 {ó}{{\'o}}1 {ú}{{\'u}}1
             {Á}{{\'A}}1 {É}{{\'E}}1 {Í}{{\'I}}1 {Ó}{{\'O}}1 {Ú}{{\'U}}1
             {ñ}{{\~n}}1 {Ñ}{{\~N}}1 {ü}{{\"u}}1 {Ü}{{\"U}}1
}

\lstset{style=sqlstyle}

\setlist[itemize]{noitemsep, topsep=0pt, leftmargin=*, itemsep=0pt}
\setlist[enumerate]{noitemsep, topsep=0pt, leftmargin=*, itemsep=0pt}
\setlength{\columnsep}{0.6cm}
\setlength{\parindent}{0pt}
\setlength{\parskip}{0.5pt}

\title{\textbf{Tema 4: PL/SQL Avanzado}}
\author{}
\date{}

\begin{document}

\tiny

\maketitle

\section*{1. Procedimientos Almacenados}

\textbf{CREATE PROCEDURE}: Define un procedimiento almacenado reutilizable.

\begin{lstlisting}
CREATE PROCEDURE nombre_proc(
  param1 IN tipo,
  param2 OUT tipo,
  param3 IN OUT tipo
) AS
  -- Declaraciones
  variable tipo;
BEGIN
  -- Código del procedimiento
  SELECT col INTO variable FROM tabla WHERE id = param1;
  param2 := variable * 2;
EXCEPTION
  WHEN NO_DATA_FOUND THEN
    param2 := 0;
END;
/
\end{lstlisting}

\textbf{Parámetros}:\\
\texttt{IN}: Parámetro de entrada (solo lectura).\\
\texttt{OUT}: Parámetro de salida (solo escritura).\\
\texttt{IN OUT}: Parámetro de entrada/salida.

\textbf{Ejecución}:
\begin{lstlisting}
EXECUTE nombre_proc(val1, :var_out, :var_inout);
-- O en bloque PL/SQL:
BEGIN
  nombre_proc(10, resultado, contador);
END;
/
\end{lstlisting}

\textbf{Eliminar procedimiento}:
\begin{lstlisting}
DROP PROCEDURE nombre_proc;
\end{lstlisting}

\section*{2. Funciones}

\textbf{CREATE FUNCTION}: Define una función que devuelve un valor.

\begin{lstlisting}
CREATE FUNCTION calc_bonus(
  salario NUMBER,
  años NUMBER
) RETURN NUMBER AS
  bonus NUMBER;
BEGIN
  IF años > 5 THEN
    bonus := salario * 0.15;
  ELSE
    bonus := salario * 0.10;
  END IF;
  RETURN bonus;
EXCEPTION
  WHEN OTHERS THEN
    RETURN 0;
END;
/
\end{lstlisting}

\textbf{Uso en consultas}:
\begin{lstlisting}
SELECT nombre, salario, 
       calc_bonus(salario, años_servicio) AS bonus
FROM empleados;
\end{lstlisting}

\textbf{Diferencias con procedimientos}:
\begin{itemize}
\item Función: RETURN obligatorio, puede usarse en SELECT
\item Procedimiento: Sin RETURN, solo ejecutable con EXECUTE/CALL
\end{itemize}

\section*{3. Triggers (Disparadores)}

\textbf{CREATE TRIGGER}: Define acciones automáticas antes/después de eventos DML.

\textbf{Sintaxis básica}:
\begin{lstlisting}
CREATE TRIGGER nombre_trigger
  {BEFORE | AFTER} {INSERT | UPDATE | DELETE}
  ON tabla
  [FOR EACH ROW]
  [WHEN (condición)]
DECLARE
  -- Declaraciones opcionales
BEGIN
  -- Código del trigger
END;
/
\end{lstlisting}

\textbf{Trigger de nivel de fila (FOR EACH ROW)}:
\begin{lstlisting}
CREATE TRIGGER audit_salario
  BEFORE UPDATE OF salario ON empleados
  FOR EACH ROW
BEGIN
  IF :NEW.salario < :OLD.salario THEN
    RAISE_APPLICATION_ERROR(-20001,
      'Salario no puede decrecer');
  END IF;
  INSERT INTO auditoria VALUES(
    :OLD.emp_id, :OLD.salario, 
    :NEW.salario, SYSDATE);
END;
/
\end{lstlisting}

\textbf{Variables :NEW y :OLD}:\\
\texttt{:NEW}: Valores nuevos (INSERT, UPDATE).\\
\texttt{:OLD}: Valores antiguos (UPDATE, DELETE).\\
No disponibles en triggers de nivel de sentencia.

\textbf{Trigger de nivel de sentencia}:
\begin{lstlisting}
CREATE TRIGGER log_cambios
  AFTER INSERT OR UPDATE OR DELETE ON tabla
BEGIN
  INSERT INTO log VALUES(USER, SYSDATE, 'DML en tabla');
END;
/
\end{lstlisting}

\textbf{INSTEAD OF trigger} (para vistas):
\begin{lstlisting}
CREATE TRIGGER upd_vista
  INSTEAD OF UPDATE ON vista_compleja
  FOR EACH ROW
BEGIN
  UPDATE tabla1 SET col = :NEW.col WHERE id = :OLD.id;
  UPDATE tabla2 SET val = :NEW.val WHERE id = :OLD.id;
END;
/
\end{lstlisting}

\textbf{Gestión de triggers}:
\begin{lstlisting}
-- Deshabilitar/habilitar
ALTER TRIGGER nombre_trigger DISABLE;
ALTER TRIGGER nombre_trigger ENABLE;
-- Eliminar
DROP TRIGGER nombre_trigger;
\end{lstlisting}

\section*{4. Paquetes (PACKAGES)}

\textbf{CREATE PACKAGE}: Agrupa procedimientos, funciones, variables y tipos relacionados.

\textbf{Especificación (interfaz pública)}:
\begin{lstlisting}
CREATE PACKAGE pkg_empleados AS
  -- Constantes públicas
  tasa_impuesto CONSTANT NUMBER := 0.21;
  -- Variables públicas
  total_procesados NUMBER := 0;
  -- Funciones públicas
  FUNCTION calc_neto(bruto NUMBER) RETURN NUMBER;
  -- Procedimientos públicos
  PROCEDURE aumentar_salario(emp_id NUMBER, pct NUMBER);
END pkg_empleados;
/
\end{lstlisting}

\textbf{Cuerpo (implementación)}:
\begin{lstlisting}
CREATE PACKAGE BODY pkg_empleados AS
  -- Variables privadas
  contador_interno NUMBER := 0;
  
  -- Función privada
  FUNCTION validar_empleado(id NUMBER) RETURN BOOLEAN AS
    v_count NUMBER;
  BEGIN
    SELECT COUNT(*) INTO v_count FROM empleados WHERE emp_id = id;
    RETURN v_count > 0;
  END;
  
  -- Implementación función pública
  FUNCTION calc_neto(bruto NUMBER) RETURN NUMBER AS
  BEGIN
    RETURN bruto * (1 - tasa_impuesto);
  END;
  
  -- Implementación procedimiento público
  PROCEDURE aumentar_salario(emp_id NUMBER, pct NUMBER) AS
  BEGIN
    IF validar_empleado(emp_id) THEN
      UPDATE empleados SET salario = salario * (1 + pct/100)
      WHERE emp_id = emp_id;
      total_procesados := total_procesados + 1;
    END IF;
  END;
END pkg_empleados;
/
\end{lstlisting}

\textbf{Uso de paquetes}:
\begin{lstlisting}
-- Acceder a elementos públicos
DECLARE
  neto NUMBER;
BEGIN
  neto := pkg_empleados.calc_neto(5000);
  pkg_empleados.aumentar_salario(101, 10);
  DBMS_OUTPUT.PUT_LINE('Total: ' || pkg_empleados.total_procesados);
END;
/
\end{lstlisting}

\section*{5. Manejo de Excepciones}

\textbf{Estructura EXCEPTION}:
\begin{lstlisting}
DECLARE
  v_salario NUMBER;
BEGIN
  SELECT salario INTO v_salario FROM empleados WHERE id = 999;
EXCEPTION
  WHEN NO_DATA_FOUND THEN
    DBMS_OUTPUT.PUT_LINE('Empleado no encontrado');
  WHEN TOO_MANY_ROWS THEN
    DBMS_OUTPUT.PUT_LINE('Múltiples empleados');
  WHEN OTHERS THEN
    DBMS_OUTPUT.PUT_LINE('Error: ' || SQLERRM);
END;
/
\end{lstlisting}

\textbf{Excepciones predefinidas}:\\
\texttt{NO\_DATA\_FOUND}: SELECT INTO sin resultados.\\
\texttt{TOO\_MANY\_ROWS}: SELECT INTO con múltiples filas.\\
\texttt{ZERO\_DIVIDE}: División por cero.\\
\texttt{DUP\_VAL\_ON\_INDEX}: Violación de clave única.\\
\texttt{INVALID\_NUMBER}: Conversión numérica inválida.\\
\texttt{VALUE\_ERROR}: Error aritmético, conversión o truncamiento.

\textbf{Excepciones definidas por usuario}:
\begin{lstlisting}
DECLARE
  salario_bajo EXCEPTION;
  v_salario NUMBER;
BEGIN
  SELECT salario INTO v_salario FROM empleados WHERE id = 101;
  IF v_salario < 1000 THEN
    RAISE salario_bajo;
  END IF;
EXCEPTION
  WHEN salario_bajo THEN
    DBMS_OUTPUT.PUT_LINE('Salario menor al mínimo');
END;
/
\end{lstlisting}

\textbf{RAISE\_APPLICATION\_ERROR}: Genera error personalizado.
\begin{lstlisting}
IF v_edad < 18 THEN
  RAISE_APPLICATION_ERROR(-20001, 
    'Edad debe ser >= 18');
END IF;
\end{lstlisting}
Rango: -20000 a -20999 para errores de usuario.

\textbf{PRAGMA EXCEPTION\_INIT}: Asocia nombre a código de error.
\begin{lstlisting}
DECLARE
  fk_violation EXCEPTION;
  PRAGMA EXCEPTION_INIT(fk_violation, -2291);
BEGIN
  DELETE FROM departamentos WHERE dept_id = 10;
EXCEPTION
  WHEN fk_violation THEN
    DBMS_OUTPUT.PUT_LINE('No se puede borrar: hay empleados');
END;
/
\end{lstlisting}

\section*{6. Cursores Avanzados}

\textbf{Cursor explícito con parámetros}:
\begin{lstlisting}
DECLARE
  CURSOR c_emp(dept_id NUMBER, min_sal NUMBER) IS
    SELECT nombre, salario FROM empleados
    WHERE departamento_id = dept_id AND salario > min_sal;
  v_nombre empleados.nombre%TYPE;
  v_salario empleados.salario%TYPE;
BEGIN
  OPEN c_emp(10, 3000);
  LOOP
    FETCH c_emp INTO v_nombre, v_salario;
    EXIT WHEN c_emp%NOTFOUND;
    DBMS_OUTPUT.PUT_LINE(v_nombre || ': ' || v_salario);
  END LOOP;
  CLOSE c_emp;
END;
/
\end{lstlisting}

\textbf{FOR...LOOP con cursor} (más simple):
\begin{lstlisting}
BEGIN
  FOR rec IN (SELECT nombre, salario FROM empleados) LOOP
    DBMS_OUTPUT.PUT_LINE(rec.nombre || ': ' || rec.salario);
  END LOOP;
END;
/
\end{lstlisting}

\textbf{Cursor con UPDATE/DELETE (FOR UPDATE)}:
\begin{lstlisting}
DECLARE
  CURSOR c_emp IS
    SELECT emp_id, salario FROM empleados
    WHERE departamento_id = 10
    FOR UPDATE OF salario NOWAIT;
BEGIN
  FOR rec IN c_emp LOOP
    UPDATE empleados SET salario = salario * 1.1
    WHERE CURRENT OF c_emp;
  END LOOP;
  COMMIT;
END;
/
\end{lstlisting}

\textbf{Atributos de cursor}:\\
\texttt{\%FOUND}: TRUE si último FETCH exitoso.\\
\texttt{\%NOTFOUND}: TRUE si último FETCH sin datos.\\
\texttt{\%ROWCOUNT}: Número de filas procesadas.\\
\texttt{\%ISOPEN}: TRUE si cursor abierto.

\textbf{Cursor REF CURSOR (dinámico)}:
\begin{lstlisting}
DECLARE
  TYPE t_cursor IS REF CURSOR;
  c_gen t_cursor;
  v_nombre VARCHAR2(100);
BEGIN
  OPEN c_gen FOR 'SELECT nombre FROM empleados WHERE dept_id = 10';
  LOOP
    FETCH c_gen INTO v_nombre;
    EXIT WHEN c_gen%NOTFOUND;
    DBMS_OUTPUT.PUT_LINE(v_nombre);
  END LOOP;
  CLOSE c_gen;
END;
/
\end{lstlisting}

\section*{7. Control Avanzado}

\textbf{CASE expresión}:
\begin{lstlisting}
v_calif := CASE
  WHEN v_nota >= 90 THEN 'Sobresaliente'
  WHEN v_nota >= 70 THEN 'Notable'
  ELSE 'Aprobado' END;
\end{lstlisting}

\textbf{Loop con etiquetas}:
\begin{lstlisting}
<<ext>> FOR i IN 1..10 LOOP
  <<int>> FOR j IN 1..10 LOOP
    EXIT ext WHEN cond;
  END LOOP int;
END LOOP ext;
\end{lstlisting}

\textbf{CONTINUE}:
\begin{lstlisting}
FOR i IN 1..100 LOOP
  CONTINUE WHEN MOD(i,2)=0;
  procesar(i);
END LOOP;
\end{lstlisting}

\section*{8. Colecciones}

\textbf{VARRAY} (array tamaño fijo):
\begin{lstlisting}
DECLARE
  TYPE t_nums IS VARRAY(10) OF NUMBER;
  v_nums t_nums := t_nums(1,2,3,4,5);
BEGIN
  v_nums.EXTEND; v_nums(6) := 100;
END;
/
\end{lstlisting}

\textbf{Nested Table} (tamaño variable):
\begin{lstlisting}
DECLARE
  TYPE t_noms IS TABLE OF VARCHAR2(50);
  v_noms t_noms := t_noms('Ana','Luis');
BEGIN
  v_noms.EXTEND(1); v_noms(3) := 'Pedro';
END;
/
\end{lstlisting}

\textbf{Associative Array} (índice string/number):
\begin{lstlisting}
DECLARE
  TYPE t_sal IS TABLE OF NUMBER INDEX BY VARCHAR2(50);
  v_sal t_sal;
BEGIN
  v_sal('Ana') := 3000; v_sal('Luis') := 4500;
END;
/
\end{lstlisting}

\textbf{Métodos}: COUNT, FIRST, LAST, NEXT(i), PRIOR(i), EXISTS(i), EXTEND(n), DELETE(i), TRIM(n).

\section*{9. SQL Dinámico}

\textbf{EXECUTE IMMEDIATE}:
\begin{lstlisting}
DECLARE
  v_sql VARCHAR2(1000); v_count NUMBER;
BEGIN
  v_sql := 'SELECT COUNT(*) FROM empleados';
  EXECUTE IMMEDIATE v_sql INTO v_count;
  v_sql := 'UPDATE empleados SET salario=salario*1.1 WHERE dept=:1';
  EXECUTE IMMEDIATE v_sql USING 10;
END;
/
\end{lstlisting}

\section*{10. Transacciones}

\textbf{AUTONOMOUS\_TRANSACTION}:
\begin{lstlisting}
CREATE PROCEDURE log_error(msg VARCHAR2) AS
  PRAGMA AUTONOMOUS_TRANSACTION;
BEGIN
  INSERT INTO error_log VALUES(SYSDATE, msg);
  COMMIT;
END;
/
\end{lstlisting}

\textbf{SAVEPOINT}:
\begin{lstlisting}
BEGIN
  UPDATE empleados SET salario = salario * 1.1;
  SAVEPOINT sp1;
  UPDATE empleados SET bonus = 500;
  ROLLBACK TO sp1;
  COMMIT;
END;
/
\end{lstlisting}

\section*{11. Ejemplos Prácticos}

\textbf{Procedimiento con cursor}:
\begin{lstlisting}
CREATE PROCEDURE ajustar_salarios(dept_id NUMBER, pct NUMBER) AS
  CURSOR c_emp IS SELECT emp_id FROM empleados
    WHERE departamento_id = dept_id FOR UPDATE;
BEGIN
  FOR rec IN c_emp LOOP
    UPDATE empleados SET salario = salario * (1 + pct/100)
    WHERE CURRENT OF c_emp;
  END LOOP;
  COMMIT;
EXCEPTION
  WHEN OTHERS THEN ROLLBACK;
END;
/
\end{lstlisting}

\textbf{Trigger de auditoría}:
\begin{lstlisting}
CREATE TRIGGER audit_salary
  AFTER UPDATE OF salario ON empleados FOR EACH ROW
BEGIN
  INSERT INTO auditoria VALUES(:OLD.emp_id, :OLD.salario, 
    :NEW.salario, USER, SYSDATE);
END;
/
\end{lstlisting}

\end{document}
